\documentclass[11pt, letterpaper]{article}
\usepackage[margin=.75in]{geometry}
\usepackage{setup}

% ==========================================
% Define parameters for this specific recital
% ==========================================
\setperformers{\emph{Junrui Liu} and \emph{Wesley Arai}}
\setdatetimeplace{Sunday, March 1, 2026}{2:00 p.m.}{Storke Tower}

\begin{document}
\begin{center}
    {\departmenthead} \\[1em]
    {\Large \recitaltype} \\[.5em]
    {\Large \textbf{\theperformers}} \\[.5em]
    {\small \thedatetimeplace} \\[1em]
\end{center}

\ruleline{program}

\vspace{1em}

\begin{programlist}
    \item
    \piece{Preludio Cou Cou}{Matthias van den Gheyn (1721--1785)}
    \editor{Ronald Barnes}

    \piece{Three Etudes for Carillon (1988)}{Stephen Rush (b. 1958)}
    \movement{I. Paced, almost hurried}
    \movement{II. Slowly, almost lost}
    \movement{III. With drive}
    \sectionperformer{Wesley Arai}

    \item
    \piece{Notule No.~1 (1955)}{Émilien Allard (1915--1976)}

    \piece{Jubilant Carol (1973)}{Roy Hamlin Johnson (1929--2020)}
    \sectionperformer{Junrui Liu}

    \item
    \piece{Introduction and Sicilienne (1982)}{Ronald Barnes (1927--1997)}

    \piece{Bell Canto (2010)}{Geert D'hollander (b. 1965)}
    \sectionperformer{Wesley Arai}

    \item
    \piece{Image 1: Brouillard (2016)}{Stefano Colletti (b. 1973)}
    \sectionperformer{Junrui Liu}

    \item
    \piece{From Classical Country Dances}{Ronald Barnes}
    \movement{I. Allegretto}
    \movement{II. Allegretto -- Grazioso}
    \sectionperformer{Junrui Liu, Primo \hspace{2em} Wesley Arai, Secondo}
\end{programlist}

\vspace{1em}

\begin{center}
    \small \it Each section will be announced by the striking of one or more bells
\end{center}

\hrulefill

\textbf{\small Welcome to Storke Tower and its Carillon!}

{\footnotesize
    The 61-bell UCSB carillon was dedicated in 1969. It was a gift of Thomas Storke, then publisher of the \textit{Santa Barbara News-Press}. The instrument, with bells cast by the Dutch bell foundry Petit \& Fritsen, is one of the finest carillons in the world with regard to bell tuning and sound quality, number of bells, and location (not near a busy street and near the ocean). The bells range in weight from about 13 pounds to 2.5 tons.
    Our instrument is one of six carillons in California, with the others being at UC Berkeley, UC Riverside, Stanford University, Christ Cathedral (formerly the Crystal Cathedral) in Garden Grove, and Trinity Cathedral in San Jose. There are over 650 carillons in the world and over 180 in North America. The carillon art is a growing one, for new instruments are installed every year. \par
}

\vfill
\calloutbox{
    \begin{center}
        \small
        \textit{Please join us for our upcoming carillon recitals:} \\[.5em]
        Sunday, March 8, 2026 at 10:00am -- UCSB Carillon Studio Student Recital \\
        Sunday, April 19, 2026 at 2:00pm -- Wesley Arai, University Carillonist \\
        Sunday, May 17, 2026 at 2:00pm -- Wesley Arai, University Carillonist
    \end{center}
}

\vfill

\newpage

\textbf{Program Notes}

\begin{small}
    \piecetitle{Preludio Cou Cou} was written by Matthias van den Gheyn, a Flemish composer best known for his carillon and organ compositions. His eleven preludes are among the few surviving original works for the carillon from the late 18th century and are still performed frequently today due to their musicality and technical virtuosity. \textit{Preludio Cou Cou} makes use of a descending third motif that is reminiscent of a cuckoo clock. Stephen Rush's \piecetitle{Three Etudes for Carillon} were commissioned by Margo Halsted (UCSB University Carillonist, 2008--2018) for the carillons at the University of Michigan. The first etude (``Paced, almost hurried'') features mixed meters with driving rhythms. The second etude (``Slowly, almost lost'') opens with a quote from the Mozart's \textit{Piano Sonata No. 11 (K. 331)}. The character of this etude is meandering and plodding, with the Mozart theme periodically reappearing. The third and final etude (``With drive'') is frenzied and, like the first etude, rhythmically intense. At times, this etude pushes the technical limits of the carillon.

    \piecetitle{Notule No.~1} is a brief work by Canadian composer Émilien Allard; the title literally translates to "a short note" or "jotting" in French. This ternary-form piece features a lyrical melody over a Siciliano rhythm in the first section. A more spirited folk dance emerges in the middle section before the work returns to the grace of the initial lyrical theme. \piecetitle{Jubilant Carol} (1973) is a piece by Roy Hamlin Johnson, who composed many significant works for the carillon over a span of more than 50 years. The melody featured throughout the piece is based in the octatonic scale, which is a series of alternating whole steps and half steps frequently used in music for the carillon due to its acoustic qualities. Johnson presents this melody in multiple textures, including in canon (with the melody starting in one voice followed by an imitation in another voice), as part of dense chords, and in the low bells with a rapid, brilliant accompanying figure in the high bells.

    Ronald Barnes played a major role in developing an American style of carillon composition. Barnes wrote countless original works, arrangers, and folk song settings for the carillon, which have become standards in the repertoire of carillonneurs in America and abroad. \piecetitle{Introduction and Sicilienne} (1982) begins with a somber yet lilting introduction, which gives way to a more cheerful, tender melody. \piecetitle{Bell Canto} was written in 2010 in celebration of the 500th anniversary of the carillon as an instrument. The title is a play on the musical term \textit{bel canto}, a style of operatic singing. The piece is subtitled ``Ravelian Waltzes in Rondo Form'' and is written in the style of French Impressionist composer Maurice Ravel, with a recurring jaunty waltz theme interspersed with a number of intermediate episodes that quote Belgian and Dutch folk songs about bells and carillons.

    \piecetitle{Image 1: Brouillard} (``Fog'') by French composer Stefano Colletti captures the essence of French Impressionism. The piece begins with restless, ambiguous harmonies that give way to a shimmering middle section dotted with broad-brush strokes of richly colored chords and glittering glissandi. Unlike the more standard "fist-style" technique of carillon playing, these sliding notes are performed with pianistic finger technique to maintain a light, glassy timbre. The work then returns to the swirling motion of the opening, bringing the fog back full circle.

    Ronald Barnes' \piecetitle{Classical Country Dances} is a collection of six short duets, each with a lively and fun air. These dances were written in 1983 during Barnes' time as University Carillonist at the University of California, Berkeley.
\end{small}

\vfill

\textbf{About the Artists}

\begin{small}

    \textbf{Junrui Liu} is a fifth-year Computer Science PhD student. His main instrument is the organ if he chances upon one to practice on, and piano otherwise. He took lessons for the former instrument in college, and has performed for recitals and services. Junrui became interested in the carillon for its fascinating sound signature and unique method of playing. When not practicing, Junrui enjoys playing badminton or watching owarai. This is Junrui's third year studying the carillon.

    \vspace{1em}

    \textbf{Wesley Arai} was appointed Lecturer and University Carillonist at the University of California, Santa Barbara in 2018. He plays the 61-bell Storke Tower carillon regularly and teaches carillon to UCSB students. Arai studied carillon with Jeff Davis as an undergraduate student at the University of California, Berkeley, where he received BA degrees in Mathematics and Statistics. While earning an MA degree in Mathematics at the University of California, Los Angeles, he continued to play the carillon and subsequently passed the Carillonneur examination of the Guild of Carillonneurs in North America. Arai then served as Associate Carillonist at the University of California, Berkeley. An active recitalist, Arai has performed extensively across the United States and abroad, most recently performing in France, Belgium, and the Netherlands. Other recent performances include recitals in Australia, the dedicatory recital for the carillon at the University of Washington, and performances at the Berkeley Carillon Festival, the Congress of the Guild of Carillonneurs in North America, the Springfield International Carillon Festival, the Perpignan (France) International Carillon Festival, and the Barcelona International Carillon Festival. Arai is also an annual recitalist at the Cathedral of St. John the Evangelist in Spokane, Washington and is a frequent recitalist at Christ Cathedral in Garden Grove, California.

\end{small}

\end{document}